\begin{statementbox}
\begin{center}\begin{tikzpicture}[font=\small,>=stealth,line join=round]
\draw[gray!40,rounded corners=3pt] (0,0) rectangle (14.2,8.3);\node at (7.1,7.700000000000001) {只看允许段：最近处最低，最远处最高};
\begin{scope}[shift={(1.8777777777777778,2.2249999999999996)},x=1.111111111111111cm,y=0.575cm]\draw[->,gray] (-0.7,0)--(4.7,0) node[right]{$x$};\draw[->,gray] (0,-1)--(0,7) node[above]{$y$};\draw[gray,thick,solid,domain=0:4,samples=180] plot(\x,{(\x-2)^2+1});\draw[blue,thick,very thick,domain=0:4,samples=180] plot(\x,{(\x-2)^2+1});\fill[orange!85!black] (0,5) circle(2.2pt);\fill[orange!85!black] (4,5) circle(2.2pt);\fill[orange!85!black] (2,1) circle(2.2pt);\draw[violet,dashed] (2,-.2)--(2,6.5);\draw[orange!85!black,very thick] (0,-.5)--(4,-.5);\node[above left] at (0,5){$(0,5)$};\node[above right] at (4,5){$(4,5)$};\node[below right] at (2,1){$(2,1)$};\end{scope}\node[anchor=west,align=left,text width=5.0cm] at (8.4,4.0) {$f=(x-2)^2+1$\\[9pt]$x\in[0,4]$\\[9pt]最近：$x=2$，最小值1\\[9pt]最远：$x=0,4$\\[9pt]最大值5，两端都取到};
\node[text width=12.6cm,align=center] at (7.1,.65) {图取a为正；开口向下时，最近与最远对应的最大最小关系反过来。};
\end{tikzpicture}\end{center}


\textbf{\textcolor{CStatement}{F030｜二次函数区间最值（轴定区间定）}}

\medskip
\textbf{\textcolor{CStatement}{这个结论解决什么问题？}}

给出一条抛物线，但只允许 $x$ 在一小段范围内变化。如何不用把每个 $x$ 都试一遍，就找到这一段图像的最高点和最低点？

\medskip
\end{statementbox}\clearpage\begin{statementbox}\textbf{\textcolor{CStatement}{使用条件}}
\begin{itemize}[leftmargin=*,itemsep=0.25em]
\item $f(x)=ax^2+bx+c$，$a\ne0$；本题中 $a,b,c$ 都是已知常数。
\item $x\in[m,n]$，$m<n$，并且两个端点 $m,n$ 都允许取到（\textbf{闭区间}）。
\item 对称轴为 $x=h$，其中 $h=-\dfrac{b}{2a}$；在本结论中，它的位置固定。
\end{itemize}

\medskip
\textbf{\textcolor{CStatement}{最重要的一句话}}

\textbf{找出区间内离对称轴最近的位置和最远的位置，就找到了最大值与最小值所在的位置。}

\begin{itemize}[leftmargin=*,itemsep=0.3em]
\item $a>0$（开口向上）：\textbf{最近处最小，最远处最大}。
\item $a<0$（开口向下）：\textbf{最近处最大，最远处最小}。
\end{itemize}

\medskip
\end{statementbox}\clearpage\begin{statementbox}\textbf{\textcolor{CStatement}{怎样找“最近处”？}}
\begin{center}\begin{tikzpicture}[font=\small,>=stealth,line join=round]
\draw[gray!40,rounded corners=3pt] (0,0) rectangle (14.2,8.4);\node at (7.1,7.800000000000001) {证明一：区间不动，先在数轴上找最近与最远};
\begin{scope}[shift={(2.5,6.1)},x=1.15cm]\draw[->,gray] (-1.6,0)--(5.8,0);\draw[blue,line width=2pt] (0,0)--(4,0);\foreach \v in {0,4}{\fill[blue] (\v,0) circle(2pt);\node[below=5pt] at (\v,0){$\v$};}\fill[violet] (-1,0) circle(2.5pt);\node[above=2pt] at (-1,0){$h=-1$};\draw[orange!85!black,line width=2pt] (0,-.16)--(0,.16);\node[anchor=west] at (-1.6,.95){轴在左侧：最近$0$；最远4};\end{scope}\begin{scope}[shift={(2.5,4.1)},x=1.15cm]\draw[->,gray] (-1.6,0)--(5.8,0);\draw[blue,line width=2pt] (0,0)--(4,0);\foreach \v in {0,4}{\fill[blue] (\v,0) circle(2pt);\node[below=5pt] at (\v,0){$\v$};}\fill[violet] (2,0) circle(2.5pt);\node[above=2pt] at (2,0){$h=2$};\draw[orange!85!black,line width=2pt] (2,-.16)--(2,.16);\node[anchor=west] at (-1.6,.95){轴在区间内：最近$2$；最远0与4};\end{scope}\begin{scope}[shift={(2.5,2.1)},x=1.15cm]\draw[->,gray] (-1.6,0)--(5.8,0);\draw[blue,line width=2pt] (0,0)--(4,0);\foreach \v in {0,4}{\fill[blue] (\v,0) circle(2pt);\node[below=5pt] at (\v,0){$\v$};}\fill[violet] (5,0) circle(2.5pt);\node[above=2pt] at (5,0){$h=5$};\draw[orange!85!black,line width=2pt] (4,-.16)--(4,.16);\node[anchor=west] at (-1.6,.95){轴在右侧：最近$4$；最远0};\end{scope}
\node[text width=12.6cm,align=center] at (7.1,.65) {蓝段为[0,4]，橙色短标记是最近位置；轴在区间内时可以直接取轴的位置。};
\end{tikzpicture}\end{center}


\[
 x_{\text{近}}=\begin{cases}
 m,&h<m,\\
 h,&m\le h\le n,\\
 n,&h>n.
 \end{cases}
\]
也就是说：顶点在区间内，就选顶点；顶点在左侧，就选左端点；顶点在右侧，就选右端点。

\medskip
\end{statementbox}\clearpage\begin{statementbox}\textbf{\textcolor{CStatement}{怎样找“最远处”？}}
\begin{center}\begin{tikzpicture}[font=\small,>=stealth,line join=round]
\draw[gray!40,rounded corners=3pt] (0,0) rectangle (14.2,8.6);\node at (7.1,8.0) {比远近：看对称轴在区间中点的哪一侧};
\begin{scope}[shift={(2,6.3)},x=1.1cm]\draw[->,gray] (-.5,0)--(8.8,0);\draw[blue,thick] (0,0)--(8,0);\foreach \v in {0,8}{\fill[blue] (\v,0) circle(2pt);\node[below=4pt] at (\v,0){$\v$};}\draw[violet,dashed] (4,-.2)--(4,.15);\fill[orange!85!black] (3,0) circle(2.5pt);\node[above=6pt] at (3,0){$h=3$};\node[anchor=west] at (0,1){左端距离3，右端距离5：右端更远};\end{scope}\begin{scope}[shift={(2,4.4)},x=1.1cm]\draw[->,gray] (-.5,0)--(8.8,0);\draw[blue,thick] (0,0)--(8,0);\foreach \v in {0,8}{\fill[blue] (\v,0) circle(2pt);\node[below=4pt] at (\v,0){$\v$};}\draw[violet,dashed] (4,-.2)--(4,.15);\fill[orange!85!black] (6,0) circle(2.5pt);\node[above=6pt] at (6,0){$h=6$};\node[anchor=west] at (0,1){左端距离6，右端距离2：左端更远};\end{scope}\begin{scope}[shift={(2,2.5)},x=1.1cm]\draw[->,gray] (-.5,0)--(8.8,0);\draw[blue,thick] (0,0)--(8,0);\foreach \v in {0,8}{\fill[blue] (\v,0) circle(2pt);\node[below=4pt] at (\v,0){$\v$};}\draw[violet,dashed] (4,-.2)--(4,.15);\fill[orange!85!black] (4,0) circle(2.5pt);\node[above=6pt] at (4,0){$h=4$};\node[anchor=west] at (0,1){左端距离4，右端距离4：两端同远};\end{scope}
\node[text width=12.6cm,align=center] at (7.1,.65) {同尺度；紫线为中点4，轴位于中点时两端才同远。};
\end{tikzpicture}\end{center}



最远的位置\textbf{一定是两个端点之一}。记区间中点 $s=\dfrac{m+n}{2}$，则
\[
\begin{array}{c|c}
\text{对称轴的位置}&\text{最远的位置}\\ \hline
h<s&n\text{（右端点）}\\
h>s&m\text{（左端点）}\\
h=s&m,n\text{ 一样远}
\end{array}
\]
这个判断与开口向上还是向下无关；开口方向只决定“最远处”取最大值还是最小值。

\medskip
\end{statementbox}\clearpage\begin{statementbox}\textbf{\textcolor{CStatement}{常规写法（便于直接套题）}}
\[
\begin{array}{ll}
 a>0:& f_{\max}=\max\{f(m),f(n)\},\\
 &f_{\min}=\begin{cases}
 f(h),&m\le h\le n,\\
 \min\{f(m),f(n)\},&h<m\text{ 或 }h>n;
 \end{cases}\\[1em]
 a<0:& f_{\min}=\min\{f(m),f(n)\},\\
 &f_{\max}=\begin{cases}
 f(h),&m\le h\le n,\\
 \max\{f(m),f(n)\},&h<m\text{ 或 }h>n.
 \end{cases}
\end{array}
\]

\textbf{边界也算在内：} $h=m$ 或 $h=n$ 时，顶点正好是端点，仍然能够取到顶点值。
\end{statementbox}
